\BourbakiExercises{Exercises}

\begin{enumerate}
\def\labelenumi{\arabic{enumi})}
\item
  Let A be a commutative ring, G a group, and H a subgroup of G of finite index. Suppose that the element \((G:H)1_{A}\) of A is invertible. Let M be an A{[}G{]}-module and N an A{[}G{]}-submodule of M. Prove that if N admits a supplementary A{[}H{]}-submodule in M, then it admits a supplementary A{[}G{]}-submodule in M (adapt the proof of Theorem 1 of VIII, p.~401). Deduce that an A{[}G{]}-module that is projective as an A{[}H{]}-module is projective.
\item
  Let A be a commutative ring and G a group.
\end{enumerate}

\begin{enumerate}
\def\labelenumi{\alph{enumi})}
\item
  Suppose that the algebra A{[}G{]} of the group G is an absolutely flat ring (VIII, p.~171, Exercise 27). Prove that the ring A is absolutely flat. Let H be a subgroup of G that admits a finite generating subset; prove that the left ideal \(I_{H}\) of A{[}G{]} (VIII, p.~22, Exercise 25) admits a supplementary in A{[}G{]}, and deduce that H is finite (observe that in the opposite case, the annihilator of \(I_{H}\) would be reduced to (0)).
\item
  Keep the assumption of a). Let g be an element of G of finite order n.~Prove that \(n1_{A}\) is invertible in A (let x be an element of A such that \((1-e_{g})x(1-e_{g})=(1-e_{g})\) ; construct an element y of A such that \(1-(1-e_{g})x=y(1+e_{g}+\cdots+e_{g^{n-1}})\) ). Deduce that the order of every finite subgroup of G is invertible in A.
\item
  Suppose that A is an absolutely flat ring and that every finite subset of G generates a finite subgroup with cardinal invertible in A. Prove that the ring A{[}G{]} is absolutely flat (reduce to the case when G is finite, and use Exercise 1).
\end{enumerate}

\begin{enumerate}
\def\labelenumi{\arabic{enumi})}
\setcounter{enumi}{2}
\tightlist
\item
  Let A be a commutative ring and G a group.
\end{enumerate}

\begin{enumerate}
\def\labelenumi{\alph{enumi})}
\item
  Suppose that the algebra A{[}G{]} of the group G is semisimple. Prove that A is semisimple, that G is finite, and that its order is invertible in A (use Exercise 2 and Exercise 25 of VIII, p.~22).
\item
  Suppose that the algebra A{[}G{]} is isomorphic to the product of a finite family of fields. Prove that every subgroup of G is normal (use a), and observe that every idempotent in A{[}G{]} is central).
\end{enumerate}

In Exercises 4 through 26 below, G is a finite group, and K is an algebraically closed field; we assume that the order of G is invertible in K.

\begin{enumerate}
\def\labelenumi{\arabic{enumi})}
\setcounter{enumi}{3}
\item
  Let \((\mathrm{M},\pi)\) be a faithful \(^{(2)}\) finite-dimensional representation of G and \((\mathrm{N},\rho)\) an irreducible representation; suppose that the character \(\chi_{\pi}\) has exactly s distinct values. Prove that there exists an integer n \textless{} s such that the K{[}G{]}-module N is isomorphic to a submodule of \(M^{\otimes n}\) (calculate \(\langle\chi_{\rho},\chi_{\pi}^{n}\rangle\) , observing that the equality \(\chi_{\pi}(g)=\dim M\) is equivalent to g=1).
\item
  Let \((\mathrm{M},\pi)\) be a representation of G.
\end{enumerate}

\begin{enumerate}
\def\labelenumi{\alph{enumi})}
\item
  Let H be a subgroup of G, S a system of representatives in G of the left cosets (mod H), and N a linear subspace of M stable under H. The representation of G in M is isomorphic to the representation induced by that of H in N if and only if we have \(M = \bigoplus_{s \in S} sN\) .
\item
  Suppose that M is the direct sum of a family \((\mathrm{M}_{i})_{i\in\mathrm{I}}\) of subspaces and that there exists a transitive action of G on I such that \(\pi(g)(\mathrm{M}_{i})\subset\mathrm{M}_{gi}\) for \(g\in G\) , \(i\in I\) . Let \(o\in I\) , and let \(G_{o}\) be its stabilizer in G. Prove that \(\pi\) is isomorphic to the representation induced by the representation of \(G_{o}\) in \(M_{o}\) . If \(\pi\) is irreducible, then the condition that the action of G on I is transitive is automatically satisfied.
\end{enumerate}

\begin{enumerate}
\def\labelenumi{\arabic{enumi})}
\setcounter{enumi}{5}
\tightlist
\item
  Let X be a finite set on which G acts. We consider the representation \(\pi\) of G in \(K^{X}\) such that \(\pi(g)e_{x}=e_{gx}\) for \(g\in G\) , \(x\in X\) (``permutation representation'' associated with X).
\end{enumerate}

\begin{enumerate}
\def\labelenumi{\alph{enumi})}
\item
  For \(g \in G\) , prove that \(\chi_{\pi}(g)\) is equal to the number of points of X fixed by g.
\item
  Prove that the multiplicity of the unit representation in \(\pi\) is equal to the number of orbits of G in X.
\item
  Assume from now on that the action of G is transitive; fix a point x of X, and denote its stabilizer by H. Prove that \(\pi\) is isomorphic to the representation induced by the unit representation of H and that \(\langle\chi_{\pi},\chi_{\pi}\rangle\) is equal to the number of orbits of H in X.
\item
  Let V be the subspace of \(K^{X}\) consisting of the vectors for which the sum of the coordinates is zero; it is stable under G, and we denote by \(\pi_{X}\) the representation of G in V deduced from \(\pi\) . Prove that \(\pi_{X}\) is irreducible if and only if the action of G is doubly transitive (I, §5, p.~143, Exercise 14).
\end{enumerate}

\begin{enumerate}
\def\labelenumi{\arabic{enumi})}
\setcounter{enumi}{6}
\item
  For any representation V of G of finite dimension, we denote by \(\varphi(\mathrm{V})\) the dimension of the subspace of V fixed by G; we extend \(\varphi\) by linearity to a Z-linear form on \(\mathrm{R}_{\mathrm{K}}(\mathrm{G})\) . Let \((\mathrm{V}_{\lambda})_{\lambda\in\widehat{\mathrm{G}}}\) be the family of irreducible representations of G, and let \(\omega\) be the element \(\sum_{\lambda\in\widehat{\mathrm{G}}}\left[\mathrm{V}_{\lambda}\right]\left[\mathrm{V}_{\lambda}^{*}\right]\) of \(\mathrm{R}_{\mathrm{K}}(\mathrm{G})\) . Prove the formula \(\varphi(\omega^{n})=\sum_{\mathrm{C}}d(\mathrm{C})^{n-1}\) for every integer \(n\geqslant0\) (apply Proposition 8 of VIII, p.~415 and the second orthogonality relation for characters).
\item
  Let H and F be subgroups of G, and let S be a subset of G such that G is the disjoint union of the double classes FsH. Let \((N,\sigma)\) be a representation of H. For \(s\in G\) , set \(H_{s}=F\cap sHs^{-1}\) , and denote by \(\sigma_{s}\) the representation of \(H_{s}\) in N defined by \(\sigma_{s}(x)=\sigma(s^{-1}xs)\) .
\end{enumerate}

\begin{enumerate}
\def\labelenumi{\alph{enumi})}
\item
  Prove that the representation \(\operatorname{Res}_{\mathrm{F}}^{\mathrm{G}}(\operatorname{Ind}_{\mathrm{H}}^{\mathrm{G}}(\sigma))\) is isomorphic to the direct sum of the representations \(\operatorname{Ind}_{\mathrm{H}_{s}}^{\mathrm{F}}(\sigma_{s})\) for \(s \in S\) . (Denote the representation \(\operatorname{Ind}_{\mathrm{H}}^{\mathrm{G}}(\sigma)\) by \((\mathrm{M}, \pi)\) . The space M can be identified with the direct sum of the \(\pi(x)\mathrm{N}\) for \(x \in G/H\) . For \(s \in S\) , let \(\mathrm{M}(s)\) be the subspace of M generated by the \(\pi(x)\mathrm{N}\) with \(x \in FsH\) . Observe that \(\mathrm{M}(s)\) is the direct sum of the subspaces \(\pi(xs)\mathrm{N}\) for \(x \in F/H_{s}\) , and deduce from Exercise 5, b) that as a K{[}F{]}-module, \(\mathrm{M}(s)\) is isomorphic to \(\operatorname{Ind}_{\mathrm{H}_{s}}^{\mathrm{F}}(\sigma_{s})\) .
\item
  Take F = H (so that we have \(H_{s} = H \cap sHs^{-1}\) ). Prove that the representation \(\operatorname{Ind}_{\mathrm{H}}^{\mathrm{G}}(\sigma)\) is irreducible if and only if \(\sigma\) is irreducible and for every \(s \in G - H\) , the supports (VIII, p.~66) of the representations \(\sigma_{s}\) and \(\operatorname{Res}_{\mathrm{H}_{s}}^{\mathrm{H}}(\sigma)\) are disjoint (``Mackey's criterion'': apply a) and Frobenius reciprocity).
\item
  Assume that H is normal; then \(\operatorname{Ind}_{\mathrm{H}}^{\mathrm{G}}(\sigma)\) is irreducible if and only if \(\sigma\) is irreducible and is not isomorphic to any of its conjugates \(\sigma \circ \operatorname{int}(g)\) for \(g \in G - H\) .
\end{enumerate}

\begin{enumerate}
\def\labelenumi{\arabic{enumi})}
\setcounter{enumi}{8}
\tightlist
\item
  Let \((\mathrm{M},\pi)\) be an irreducible representation of G and H a normal subgroup of G. a) Let \(M=\bigoplus_{i\in I}M_{i}\) be the decomposition of the K{[}H{]}-module M as a direct sum of isotypical submodules. Prove that the action of G on M permutes the submodules \(M_{i}\) transitively. Deduce that \((\mathrm{M},\pi)\) is isomorphic to the representation induced by a subgroup of G containing H (apply Exercise 5, b)).
\end{enumerate}

\begin{enumerate}
\def\labelenumi{\alph{enumi})}
\setcounter{enumi}{1}
\tightlist
\item
  Let \(\mathcal{S}\) be the support of the K{[}H{]}-module M; there exists an integer e such that M is K{[}H{]}-isomorphic to \(\sum_{\lambda\in \mathcal{S}}\lambda^{e}\) .
\end{enumerate}

\begin{enumerate}
\def\labelenumi{\arabic{enumi})}
\setcounter{enumi}{9}
\tightlist
\item
  Let \((\mathrm{M},\pi)\) be an irreducible representation of G and H a normal subgroup of G. a) Suppose that H is equal to the center of G. Prove that the degree of \(\pi\) divides the index of H in G. (For \(m \geqslant 1\) , let \(H_{m}\) be the subgroup of \(H^{m}\) consisting of the elements \((h_{1},\ldots,h_{m})\) such that \(h_{1}\ldots h_{m}=1\) . Observe that the representation \(\pi^{\times m}:G^{m}\to\operatorname{End}_{K}(M^{\otimes m})\) defines an irreducible representation of \(G^{m}/H_{m}\) , and apply Corollary 2 of Proposition 10 (VIII, p.~420) to it.)
\end{enumerate}

\begin{enumerate}
\def\labelenumi{\alph{enumi})}
\setcounter{enumi}{1}
\tightlist
\item
  Suppose that H is commutative. Prove that the degree of \(\pi\) divides (G : H). (Using Exercise 9, a), and reasoning by induction on the cardinal of G, reduce to the case when the restriction of \(\pi\) to H is isotypical. Prove that \(\pi(\mathrm{H})\) is central in \(\pi(\mathrm{G})\) , and apply a).
\end{enumerate}

\begin{enumerate}
\def\labelenumi{\arabic{enumi})}
\setcounter{enumi}{10}
\tightlist
\item
  Let A and H be subgroups of G. Suppose that A is commutative and normal and that G is the semidirect product of H and A. The group H acts on \(\widehat{\mathrm{A}} = \mathrm{Hom}(\mathrm{A}, \mathrm{K}^{*})\) ; choose a system of representatives \((\alpha_{i})_{i \in \widehat{\mathrm{A}}/\mathrm{H}}\) of the orbits of H in \(\widehat{A}\) . For \(i \in \widehat{A}/H\) , let \(H_{i}\) be the stabilizer of \(\alpha_{i}\) in H, and let \(G_{i} = AH_{i}\) . For every irreducible representation \(\sigma\) of \(H_{i}\) , define a representation \(\rho_{i,\sigma}\) of \(G_{i}\) by setting \(\rho_{i,\sigma}(ah) = \alpha_{i}(a)\sigma(h)\) for \(a \in A\) and \(h \in H\) . Prove that the representations \(\operatorname{Ind}_{G_{i}}^{G}(\rho_{i,\sigma})\) , for \(i \in \widehat{A}/H\) and \(\sigma \in \widehat{H}_{i}\) , are irreducible and pairwise nonisomorphic, and that we obtain all irreducible representations of G this way (apply Exercise 9).
\end{enumerate}

¶ 12) Let H be a subgroup of G; suppose that for every \(g \in G - H\) , we have the equality \(H \cap gHg^{-1} = \{1\}\) .

\begin{enumerate}
\def\labelenumi{\alph{enumi})}
\item
  Let u be a central function on H. Prove that the function \(\operatorname{Ind}_{\mathrm{H}}^{\mathrm{G}}(u)\) coincides with u on \(H - \{1\}\) . If v is another central function on H such that \(v(1) = 0\) , then we have \(\langle \operatorname{Ind}_{\mathrm{H}}^{\mathrm{G}}(u), \operatorname{Ind}_{\mathrm{H}}^{\mathrm{G}}(v) \rangle_{\mathrm{G}} = \langle u, v \rangle_{\mathrm{H}}\) .
\item
  Let \(\lambda \in \widehat{\mathrm{H}}\) ; we denote by \(\chi_{\lambda}^{\prime}\) the function \(\chi_{\lambda} - d_{\lambda}\) on H and by \(\theta_{\lambda}\) the function \(d_{\lambda} + \mathrm{Ind}_{\mathrm{H}}^{\mathrm{G}}(\chi_{\lambda}^{\prime})\) on G. Prove that \(\theta_{\lambda}\) is the character of an irreducible representation of G of degree \(d_{\lambda}\) whose restriction to H is isomorphic to \(\lambda\) (use a) to calculate \(\langle\theta_{\lambda},\theta_{\lambda}\rangle\) , \(\langle\theta_{\lambda},1\rangle\) , and \(\langle\mathrm{Res}_{\mathrm{H}}^{\mathrm{G}}(\theta_{\lambda}),\chi_{\lambda}\rangle\) ).
\item
  Set \(\theta = \sum_{\lambda} d_{\lambda} \theta_{\lambda}\) . Prove that we have \(\theta(h) = 0\) if \(h \in \mathrm{H} - \{1\}\) and \(\theta(g) = \theta(1) = \mathrm{Card}(\mathrm{H})\) if \(g\) does not belong to any conjugates of \(\mathrm{H}\) .
\item
  Let \(L'\) be the set of elements of G that do not belong to any conjugates of H, and set \(L = L' \cup \{1\}\) . Prove that L is a normal subgroup of G (deduce from c) that L is the kernel of a representation with character \(\theta\) .
\item
  Prove that G is the semidirect product of H and L (calculate the order of L). f) Let \(h \in H - \{1\}\) . Prove that the automorphism \(\operatorname{int}(h)\) restricted to L has no fixed point other than 1. Deduce that the order of H divides \(\operatorname{Card}(L) - 1\) .
\end{enumerate}

\(g\) ) Prove that if the subgroup H has even order, then L is commutative (apply Exercise 23, \(d\) ) of I, §6, p.~145).

\begin{enumerate}
\def\labelenumi{\arabic{enumi})}
\setcounter{enumi}{12}
\tightlist
\item
  Let \(\mathcal{H}\) be a set of subgroups of \(G\) .
\end{enumerate}

\begin{enumerate}
\def\labelenumi{\alph{enumi})}
\tightlist
\item
  Prove that the following properties are equivalent:
\end{enumerate}

\begin{enumerate}
\def\labelenumi{(\roman{enumi})}
\item
  The union of the conjugates of the subgroups belonging to \(\mathcal{H}\) is equal to G.
\item
  Every character of \(G\) is a linear combination with rational coefficients of characters induced by the characters of subgroups belonging to \(\mathcal{H}\) .
\end{enumerate}

(Deduce from Proposition 6 of VIII, p.~412 that property (ii) is equivalent to the injectivity of the restriction homomorphism \(\mathcal{Z}_{\mathrm{K}}(\mathrm{G})\to \bigoplus_{\mathrm{H}\in \mathcal{H}}\mathcal{Z}_{\mathrm{K}}(\mathrm{H}).\))

\begin{enumerate}
\def\labelenumi{\alph{enumi})}
\setcounter{enumi}{1}
\item
  Let \(\mathcal{C}\) be the set of cyclic subgroups of G; it obviously has property (i). For every cyclic subgroup C of G, we denote by \(\varphi_{\mathrm{C}}\) the function on C with value \(\operatorname{Card}(\mathrm{C})\) on the generators of C and 0 on the other elements. Prove the equality \(\operatorname{Card}(\mathrm{G}) = \sum_{\mathrm{C} \in \mathcal{C}} \operatorname{Ind}_{\mathrm{C}}^{\mathrm{G}}(\varphi_{\mathrm{C}})\) .
\item
  Let \(C \in \mathcal{C}\) . Prove that the function \(\varphi_{C}\) belongs to the subgroup \(\mathrm{R}_{\mathrm{K}}(\mathrm{G})\) of \(\mathcal{Z}_{\mathrm{K}}(\mathrm{G})\) generated by the characters (reason by induction on the order of C, using b)).
\end{enumerate}

\(d\) ) Let \(\chi\) be the character of an irreducible representation G. Prove the equality \(\chi = \mathrm{Card}(\mathrm{G})^{-1}\sum_{\mathrm{C}\in \mathcal{C}}\mathrm{Ind}_{\mathrm{C}}^{\mathrm{G}}(\varphi_{\mathrm{C}}\mathrm{Res}_{\mathrm{C}}^{\mathrm{G}}(\chi))\) , which gives an explicit formula for property (ii) of \(a\) ).

\begin{enumerate}
\def\labelenumi{\arabic{enumi})}
\setcounter{enumi}{13}
\tightlist
\item
  A representation of G is called monomial if it is the direct sum of representations induced by degree 1 representations of subgroups of G.
\end{enumerate}

\begin{enumerate}
\def\labelenumi{\alph{enumi})}
\item
  Suppose that the representations of the proper subgroups of G are monomial and that G contains a noncentral commutative normal subgroup H. Prove that every faithful irreducible representation of G is monomial (observe that the restriction to H of such a representation is not isotypical, and apply Exercise 9, a)).
\item
  Deduce that all representations of supersolvable groups (I, §6, p.~146, Exercise 26), and in particular of nilpotent groups, are monomial.
\item
  Take G to be the group \(\mathbf{SL}(2,\mathbf{F}_{3})\) . Its center Z has order 2, and the group G/Z is isomorphic to the alternating group \(\mathfrak{A}_{4}\) (II, §10, p.~422, Exercise 14, g)); in particular, G is solvable. Prove that its derived group has index 3 (cf.~I, §5, p.~142,
\end{enumerate}

Exercise 10). Using formula (8) of VIII, p.~407, deduce that G admits an irreducible representation of degree 2, and prove that this representation is not monomial.

¶ 15) Suppose that all representations of G are monomial; let us prove that G is solvable. Reasoning by induction on the order of G, we may assume that every quotient group of G different from G is solvable.

\begin{enumerate}
\def\labelenumi{\alph{enumi})}
\tightlist
\item
  If G contains two distinct minimal nontrivial normal subgroups \(H_{1}\) and \(H_{2}\) , then it is solvable (consider the canonical homomorphism from G to \(G/H_{1} \times G/H_{2}\) ).
\end{enumerate}

Assume from now on that G admits a unique minimal nontrivial normal subgroup H.

\begin{enumerate}
\def\labelenumi{\alph{enumi})}
\setcounter{enumi}{1}
\tightlist
\item
  Prove that a representation of G is faithful if its restriction to H is nontrivial. c) Let \(\pi\) be a faithful representation of minimal degree; there exist subgroups \(H_{i}\) of G and representations \(\sigma_{i}\) of \(H_{i}\) of degree 1 such that \(\pi\) is isomorphic to \(\bigoplus\operatorname{Ind}_{\operatorname{H}_{i}}^{\operatorname{G}}(\sigma_{i})\) . Let \(\pi_{0}=\bigoplus\operatorname{Ind}_{\operatorname{H}_{i}}^{\operatorname{G}}(\varepsilon_{\operatorname{H}_{i}})\) , where \(\varepsilon_{H_{i}}\) denotes the unit representation. Prove that the kernel of \(\pi_{0}\) is commutative and nontrivial (observe that \(\pi_{0}\) is irreducible); deduce that G is solvable.
\end{enumerate}

\begin{enumerate}
\def\labelenumi{\arabic{enumi})}
\setcounter{enumi}{15}
\tightlist
\item
  Suppose that the characteristic of K is zero. We say that an element x of K is integral over Z if it belongs to a subring of K that is a finitely generated Z-module or, equivalently, if the subring \(Z[x]\) of K is a finitely generated Z-module \(^{*}\) (cf.~Comm. Alg., V, §1, No.~1, p.~303).
\end{enumerate}

\begin{enumerate}
\def\labelenumi{\alph{enumi})}
\item
  Prove that the set of elements of K integral over Z is a subring of K (for x, y integral over Z, consider the subring \(Z[x, y]\) ).
\item
  An element \(x\) of \(K\) is integral over \(\mathbf{Z}\) if and only if it is a root of a monic polynomial in \(\mathbf{Z}[X]\) (when \(x\) is integral, adapt the proof of Lemma 1 of VIII, p.~81 to show that it satisfies an equation of the form \(\det(xI_m - A) = 0\) with \(A \in \mathbf{M}_m(\mathbf{Z})\) ).
\item
  If x is integral over Z, then the same holds for its conjugates \(x_{1},\ldots,x_{m}\) over Q; the element \(x_{1}\cdots x_{m}\) belongs to Z (use the lemma of VIII, p.~416).
\end{enumerate}

\begin{enumerate}
\def\labelenumi{\arabic{enumi})}
\setcounter{enumi}{16}
\tightlist
\item
  Let \(\mathrm{C}\) be a conjugacy class of \(\mathrm{G}\) and \(\pi\) an irreducible representation of \(\mathrm{G}\) of degree \(d\) .
\end{enumerate}

\begin{enumerate}
\def\labelenumi{\alph{enumi})}
\item
  Prove that the element \(d^{-1}\operatorname{Card}(\mathrm{C})\chi(\mathrm{C})\) of K is integral over Z (reason as in the proof of Proposition 9 of VIII, p.~415).
\item
  Suppose that d and \(\operatorname{Card}(\mathbf{C})\) are mutually prime integers. Prove that the element \(d^{-1}\chi(\mathbf{C})\) is integral over Z.
\end{enumerate}

* \(c\) ) Suppose, moreover, \(\mathrm{K} = \mathbf{C}\) . Prove that all conjugates of \(d^{-1}\chi (\mathrm{C})\) over \(\mathbf{Q}\) have absolute value \(\leqslant 1\) . Deduce from Exercise 16, \(c\) ) that \(|\chi (\mathrm{C})|\) is zero or equal to \(d\) . Conclude that if \(\chi (\mathrm{C})\) is not zero, then \(\pi (x)\) is a homothety for every \(x\in \mathrm{C}\) . In particular, if \(\pi\) is faithful and C is not reduced to one element, then we have \(\chi (\mathrm{C}) = 0.\ast\)

\begin{enumerate}
\def\labelenumi{\arabic{enumi})}
\setcounter{enumi}{17}
\tightlist
\item
  Suppose that the group G is simple but not cyclic.
\end{enumerate}

\begin{enumerate}
\def\labelenumi{\alph{enumi})}
\item
  Prove that there is no conjugacy class different from \(\{1\}\) whose cardinal is a power of a prime number. (Let C be a class of cardinal \(p^{r}\) with \(r \geqslant 1\) . Deduce from Exercise 17, b) that \(p^{-1}\chi_{\lambda}(1)\chi_{\lambda}(C)\) is integral over Z when \(\lambda\) is different from the unit representation, and deduce a contradiction using the second orthogonality relation for characters).
\item
  Prove that the order of G is divisible by at least three distinct prime numbers (apply a) to the conjugacy class of a central element \(x \neq 1\) in a Sylow subgroup of G).
\item
  Prove that a finite group whose order is divisible by at most two distinct prime numbers is solvable (``Burnside's theorem'': reason by induction on the cardinal of the group, using b)).
\end{enumerate}

*19) Let \(\pi\) be an irreducible complex representation of G of degree \(>1\) . Prove that there exists an element \(x\) of G such that \(\chi_{\pi}(x)=0\) . (Let \(\{\lambda_1,\ldots,\lambda_s\}\) be the orbit of \(\pi\) in \(\widehat{G}\) for the action of the group \(\operatorname{Gal}(\mathbf{C}/\mathbf{Q})\) , and let \(\chi_1,\ldots,\chi_s\) be the corresponding characters. Let \(g\in G\) ; if \(\chi(g)\neq 0\) , then prove as in Exercise 17, \(c\) ) that we have \(|\chi_1(g)\ldots\chi_s(g)|\geqslant 1\) , and therefore \(\sum_i|\chi_i(g)|^2\geqslant s\) (FRV, III, §1, No.~1, p.~93, Proposition 2). If this holds for every \(g\in G\) , then deduce from the orthogonality relation for characters that we have equality, which leads to a contradiction by taking \(g=1\) .)*

\begin{enumerate}
\def\labelenumi{\arabic{enumi})}
\setcounter{enumi}{19}
\tightlist
\item
  Let G be a finite group, H a subgroup of G, \(\sigma\) a finite-dimensional representation of H, and \(\rho\) the representation of G induced by \(\sigma\) . Let C be a conjugacy class of G; the set \(C \cap H\) is the disjoint union of a finite family \((D_{i})_{i \in I}\) of conjugacy classes of H.
\end{enumerate}

\begin{enumerate}
\def\labelenumi{\alph{enumi})}
\item
  Prove the formula \(\chi_{\rho}(\mathrm{C}) = \frac{(\mathrm{G:H})}{\mathrm{Card}(\mathrm{C})}\sum_{i\in \mathrm{I}}\mathrm{Card}(\mathrm{D}_i)\chi_\sigma (\mathrm{D}_i)\) .
\item
  If \(\sigma\) is the trivial representation, then \(\chi_{\rho}(\mathrm{C})=(\mathrm{G}:\mathrm{H})\operatorname{Card}(\mathrm{C})^{-1}\operatorname{Card}(\mathrm{C}\cap\mathrm{H})\) holds.
\end{enumerate}

¶ 21) Let \(n \in \mathbf{N}\) . Denote by \(\mathcal{D}_{n}\) the subset of \(\mathbf{N}^{n}\) consisting of the elements \((p_{1},\ldots,p_{n})\) such that \(p_{1} \geqslant \cdots \geqslant p_{n}\) ; endow it with the lexicographical order. For \(\mathbf{p} \in \mathcal{D}_{n}\) , denote by \(\mathbf{p}'\) the element of \(\mathcal{D}_{n}\) defined by \(p_{i}' = p_{i} + n - i\) . Let \(\mathbf{X} = (\mathrm{X}_{1},\ldots,\mathrm{X}_{n})\) be a sequence of variables; set \(\Delta(\mathbf{X}) = \prod_{i<j}(\mathrm{X}_{i} - \mathrm{X}_{j})\) .

\begin{enumerate}
\def\labelenumi{\alph{enumi})}
\item
  Let \(\mathrm{P} \in \mathbf{Z}[\mathbf{X}]\) be an antisymmetric polynomial, that is, a polynomial satisfying the relation \(\mathrm{P}(\mathrm{X}_{\sigma(1)}, \ldots, \mathrm{X}_{\sigma(n)}) = \varepsilon(\sigma)\mathrm{P}(\mathrm{X}_1, \ldots, \mathrm{X}_n)\) for every element \(\sigma \in \mathfrak{S}_n\) . Prove that there exists a symmetric polynomial \(\mathrm{Q} \in \mathbf{Z}[\mathbf{X}]\) such that \(\mathrm{P} = \Delta\mathrm{Q}\) .
\item
  For every \(\mathbf{p} \in \mathcal{D}_n\) , set \(S_{\mathbf{p}}(\mathbf{X}) = \Delta(\mathbf{X})^{-1} \det(X_j^{p_i'})_{1 \leqslant i,j \leqslant n}\) . Prove that \(S_{\mathbf{p}}\) is a homogeneous symmetric polynomial of degree \(\sum p_i\) with integer coefficients.
\item
  If \(\mathbf{q} = (q_i)_{1\leqslant i\leqslant n}\) is an element of \(\mathcal{D}_n\) and \(P\) is an element of \(\mathbf{Q}[\mathbf{X}]\) , then we denote by \([P]_{\mathbf{q}}\) the coefficient of \(\mathbf{x}^{\mathbf{q}}\) in \(P\) . Set \(K_{\mathbf{pq}} = [S_{\mathbf{p}}]_{\mathbf{q}}\) . Prove that we have \(K_{\mathbf{pp}} = 1\) and \(K_{\mathbf{pq}} = 0\) for \(\mathbf{q} > \mathbf{p}\) (expand the equality \(\det(X_j^{p_i'}) = S_{\mathbf{p}}(\mathbf{X})\Delta(\mathbf{X})\) in the lexicographical order). Deduce that the polynomials \(S_{p}\) for \(p \in \mathcal{D}_{n}\) form a basis of the Z-submodule of Z{[}X{]} consisting of the symmetric polynomials.
\item
  Let \(\mathbf{Y} = (\mathrm{Y}_1, \ldots, \mathrm{Y}_n)\) be a sequence of variables. Prove the equality of formal power series
\end{enumerate}

\[
\prod_ {1 \leqslant i, j \leqslant n} (1 - \mathrm{X} _ {i} \mathrm{Y} _ {j}) ^ {- 1} = \sum_ {\mathbf {p} \in \mathcal {D} _ {n}} \mathrm{S} _ {\mathbf {p}} (\mathbf {X})   \mathrm{S} _ {\mathbf {p}} (\mathbf {Y})
\]

(use the equality \(\det\big((1 - X_iY_j)^{-1}\big)_{1\leqslant i,j\leqslant n} = \Delta(\mathbf{X})\Delta(\mathbf{Y})\prod_{1\leqslant i,j\leqslant n}(1 - X_iY_j)^{-1}\) (cf.~III, §8, p.~639, Exercise 5) by calculating the left-hand side using the expansion of \((1 - T)^{-1}\) as a formal power series).

\begin{enumerate}
\def\labelenumi{\alph{enumi})}
\setcounter{enumi}{4}
\tightlist
\item
  Let \(\mathrm{P} \in \mathbf{Q}[\mathbf{X}]\) be a symmetric polynomial. For every \(\mathbf{p} \in \mathcal{D}_n\) , set \(\omega_{\mathbf{p}}(\mathrm{P}) = [\mathrm{P}\Delta]_{\mathbf{p}'}\) . Prove the equality \(\omega_{\mathbf{p}}(\mathrm{S}_{\mathbf{q}}) = \delta_{\mathbf{pq}}\) . Deduce that for every \(\mathbf{q} \in \mathcal{D}_n\) , we have \([\mathrm{P}]_{\mathbf{q}} = \sum_{\mathbf{p} \in \mathcal{D}_n} \mathrm{K}_{\mathbf{pq}} \omega_{\mathbf{p}}(\mathrm{P})\) .
\end{enumerate}

\(f)\) For \(\boldsymbol{\alpha}\in\mathbf{N}^{n}\) , set \(f^{\boldsymbol{\alpha}}(\mathbf{X})=\prod_{j=1}^{n}(\sum_{i}\mathrm{X}_{i}^{j})^{\alpha_{j}}\) and \(c_{\boldsymbol{\alpha}}=\prod_{j=1}^{n}\frac{1}{j^{\alpha_{j}}\alpha_{j}!}\) . Prove that we have \(\sum_{\boldsymbol{\alpha}\in\mathbf{N}^{n}}c_{\boldsymbol{\alpha}}\omega_{\mathbf{p}}(f^{\boldsymbol{\alpha}})\omega_{\mathbf{q}}(f^{\boldsymbol{\alpha}})=\delta_{\mathbf{pq}}\) for \(\mathbf{p},\mathbf{q}\) in \(\mathcal{D}_{n}\) (observe that we have

\[
\prod_ {1 \leqslant i, j \leqslant n} (1 - \mathrm{X} _ {i} \mathrm{Y} _ {j}) ^ {- 1} = \prod_ {k = 0} ^ {\infty} \exp \bigl (\frac {1}{k} (\sum_ {i} \mathrm{X} _ {i} ^ {k}) (\sum_ {j} \mathrm{Y} _ {j} ^ {k}) \bigr) = \sum_ {\boldsymbol {\alpha} \in \mathbf {N} ^ {n}} c _ {\boldsymbol {\alpha}} f ^ {\boldsymbol {\alpha}} (\mathbf {X}) f ^ {\boldsymbol {\alpha}} (\mathbf {Y}) \bigr).
\]

¶ 22) Let n be a natural number. Denote by \(\mathcal{S}_{n}\) the set of decreasing finite sequences of strictly positive integers with sum n.~Let \(\mathbf{p} = (p_{i})_{1 \leqslant i \leqslant d}\) be an element of \(\mathcal{S}_{n}\) . For \(1 \leqslant i \leqslant d\) , denote by \(P_{i}\) the set of integers r such that \(p_{1} + \cdots + p_{i-1} < r \leqslant p_{1} + \cdots + p_{i}\) , and for \(1 \leqslant j \leqslant p_{1}\) , denote by \(P_{j}'\) the set of integers of the form \(p_{1} + \cdots + p_{i-1} + j\) with \(1 \leqslant i \leqslant d\) and \(p_{i} \geqslant j\) . The subsets \(P_{i}\) , on the one hand, and the subsets \(P_{j}'\) , on the other, form a partition of the interval {[}1, n{]} in N.

(We often represent this situation by a diagram, called a Young diagram: in the example below, we have n = 9 and \(\mathbf{p} = (4, 2, 2, 1)\) . The set \(P_{i}\) consists of the elements of the i-th row, and \(P_{j}^{\prime}\) of those of the j-th column.)

1

2

3

4

5

6

7

8

9

Let \(\mathcal{P}\) (resp. \(\mathcal{P}'\) ) be the subgroup of \(\mathfrak{S}_n\) consisting of the permutations \(\sigma\) such that \(\sigma(\mathrm{P}_i) \subset \mathrm{P}_i\) for every \(i\) (resp. \(\sigma(\mathrm{P}_j') \subset \mathrm{P}_j'\) for every \(j\) ).

\begin{enumerate}
\def\labelenumi{\alph{enumi})}
\item
  Prove that we have \(\mathcal{P} \cap \mathcal{P}' = \{1\}\) .
\item
  Let \(\sigma \in \mathfrak{S}_n\) . Show that if \(\mathrm{Card}(\mathrm{P}_i \cap \sigma(\mathrm{P}_j')) \leqslant 1\) for every \(i\) and \(j\) , then \(\sigma\) belongs to \(\mathcal{PP}'\) . Deduce that if \(\sigma \notin \mathcal{PP}'\) , then the subgroup \(\mathcal{P} \cap \sigma \mathcal{P}'\sigma^{-1}\) contains a transposition.
\item
  In the group algebra \(\mathbf{Z}[\mathfrak{S}_n]\) , set
\end{enumerate}

\[
a _ {\mathbf {p}} = \sum_ {\sigma \in \mathcal {P}} e _ {\sigma}, b _ {\mathbf {p}} = \sum_ {\tau \in \mathcal {P} ^ {\prime}} \varepsilon (\tau) e _ {\tau}, c _ {\mathbf {p}} = a _ {\mathbf {p}} b _ {\mathbf {p}}.
\]

Prove that \(c_{p}\) satisfies \(e_{\sigma}c_{p}e_{\sigma'}=\varepsilon(\sigma')c_{p}\) for every \(\sigma\in \mathcal{P}\) and \(\sigma'\in \mathcal{P}'\) and that every element that has this property belongs to \(Zc_{p}\) (write such a element as \(\sum n_{\sigma}e_{\sigma}\) , and deduce from b) that \(n_{\sigma}\) is zero for \(\sigma\notin \mathcal{P}\mathcal{P}'\) ). Deduce that for every element x of \(Z[\mathfrak{S}_n]\) , the element \(c_{p}xc_{p}\) belongs to \(Zc_{p}\) .

\(d\) ) Denote by A the \(\mathbf{Q}\) -algebra \(\mathbf{Q}[\mathfrak{S}_n]\) . Prove that the ideal \(V_{\mathbf{p}} = Ac_{\mathbf{p}}\) of A is an absolutely simple A-module. (Observe that we have \(c_{\mathbf{p}}V_{\mathbf{p}} \subset \mathbf{Q}c_{\mathbf{p}}\) by \(c\) ). Prove that if \(\mathfrak{a}\) is a left ideal strictly contained in \(V_{\mathbf{p}}\) , then we have \(c_{\mathbf{p}}\mathfrak{a} = 0\) , and therefore \(\mathfrak{a}^2 = 0\) and finally \(\mathfrak{a} = 0\) .)

\begin{enumerate}
\def\labelenumi{\alph{enumi})}
\setcounter{enumi}{4}
\tightlist
\item
  Prove that we have \(c_{\mathbf{p}}^{2} = (n!/[V_{\mathbf{p}} : \mathbf{Q}]) c_{\mathbf{p}}\) (calculate the trace of the endomorphism \(x \mapsto xc_{p}\) of A).
\end{enumerate}

¶ 23) Keep the notation of the previous exercise. Let \(\mathbf{p} = (p_i)\) and \(\mathbf{q} = (q_j)\) be two elements of \(S_n\) .

\begin{enumerate}
\def\labelenumi{\alph{enumi})}
\item
  Suppose that we have \(\mathbf{p} > \mathbf{q}\) for the lexicographical order. Prove that we have \(a_{\mathbf{p}}xb_{\mathbf{q}} = 0\) for every \(x \in \mathrm{A}\) . (Reduce to the case \(x = 1\) . Observe that if \((\mathrm{P}_i), (\mathrm{P}_j')\) denote the partitions associated with \(\mathbf{p}\) and \((\mathrm{Q}_k), (\mathrm{Q}_l')\) those associated with \(\mathbf{q}\) , then, denoting by \(s\) the smallest integer \(n\) such that \(p_n \neq q_n\) , we have \(\operatorname{Card}(\mathrm{P}_s \cap \mathrm{Q}_1') \geqslant 2\) , and construct a transposition \(\tau\) such that \(a_{\mathbf{p}}\tau = a_{\mathbf{p}}\) and \(\tau b_{\mathbf{q}} = -b_{\mathbf{q}}\) .)
\item
  Deduce that if \(p \neq q\) , then the A-modules \(V_{p}\) and \(V_{q}\) are not isomorphic (apply a), using the antiautomorphism of A that sends \(e_{\sigma}\) to \(e_{\sigma^{-1}}\) .
\item
  Prove that the mapping that sends a sequence \(\mathbf{p}\) to the representation \((\mathrm{V}_{\mathbf{p}})_{(\mathrm{K})}\) defines a bijection from \(\mathcal{S}_n\) to the set \(\mathcal{S}_{\mathrm{K}}(\mathfrak{S}_n)\) of classes of simple \(\mathrm{K}[\mathfrak{S}_n]\) -modules. d) For \(\mathbf{p} \in \mathcal{S}_n\) , denote by \(\overline{\mathbf{p}}\) the sequence defined by \(\overline{p}_j = \mathrm{Card}(\mathrm{P}_j')\) (Exercise 22). Prove that the mapping \(\mathbf{p} \mapsto \overline{\mathbf{p}}\) is an involution of \(\mathcal{S}_n\) ; the A-module \(\mathrm{V}_{\overline{\mathbf{p}}}\) is isomorphic to \(\mathrm{V}_{\mathbf{p}} \otimes_{\mathbf{Q}} \mathbf{Q}_{\varepsilon}\) , where \(\mathbf{Q}_{\varepsilon}\) denotes the 1-dimensional representation over \(\mathbf{Q}\) associated with the signature.
\item
  Describe the representations corresponding to the sequences \((1, \ldots, 1)\) and \((n)\) , as well as to the sequence \((n - 1, 1)\) .
\end{enumerate}

¶ 24) Keep the notation of Exercises 21 through 23. Let \(\mathbf{p} = (p_i)_{1 \leqslant i \leqslant d} \in \mathcal{S}_n\) ; denote by \(\mathbf{p}\) the element \((p_1, \ldots, p_d, 0, \ldots, 0)\) of \(\mathcal{D}_n\) . If \(\boldsymbol{\alpha} = (\alpha_1, \ldots, \alpha_n)\) is an element of \(\mathbf{N}^n\) such that \(\sum_i i\alpha_i = n\) , then we denote by \(C_{\boldsymbol{\alpha}}\) the conjugacy class of \(\mathfrak{S}_n\) consisting of the elements that are the product of a family of cycles with disjoint support consisting of \(\alpha_2\) cycles of order 2, \(\alpha_3\) cycles of order 3, etc.

\begin{enumerate}
\def\labelenumi{\alph{enumi})}
\item
  Prove that the representation \(\rho_{\mathbf{p}}\) of \(\mathfrak{S}_n\) in the A-module \(\mathrm{Aa}_{\mathbf{p}}\) is the representation induced by the trivial representation of \(\mathcal{P}\) . Use Exercise 20 to deduce the formula \(\chi_{\rho_{\mathbf{p}}}(\mathrm{C}_{\alpha}) = [f^{\alpha}]_{\mathbf{p}}\) .
\item
  For any \(\mathbf{q} \in \mathcal{S}_n\) , denote by \(\lambda_{\mathbf{q}}\) the representation of \(\mathfrak{S}_n\) in the A-module \(V_{\mathbf{q}}\) . Prove that there exist positive integers \(n_{\mathbf{pq}}\) such that we have \(\chi_{\rho_{\mathbf{p}}} = \sum_{\mathbf{q} \in \mathcal{S}_n} n_{\mathbf{pq}} \chi_{\lambda_{\mathbf{q}}}\) with \(n_{\mathbf{pp}} > 0\) .
\item
  Let \(\eta_{\mathbf{p}}\) be the central function on \(\mathfrak{S}_n\) with value \(\omega_{\mathbf{p}}(f^{\alpha})\) on the elements of the class \(\mathrm{C}_{\alpha}\) . Prove that \(\eta_{\mathbf{p}}\) is a linear combination of the characters \(\chi_{\lambda_{\mathbf{q}}}\) for \(\mathbf{q} \in \mathcal{S}_n\) with integer coefficients (use Exercise 21, c) and e)). Deduce that \(\eta_{\mathbf{p}}\) is equal, up to a sign, to one of the characters \(\chi_{\lambda_{\mathbf{q}}}\) . (Calculate \(\langle \eta_{\mathbf{p}}, \eta_{\mathbf{p}} \rangle\) using Exercise 21, f). d) Deduce the equality \(\chi_{\lambda_{\mathbf{p}}}(\mathrm{C}_{\alpha}) = [\Delta f^{\alpha}]_{\mathbf{p}'}\) (``Frobenius formula'').
\item
  Prove that the multiplicity of the irreducible representation \(\lambda_{p}\) in the representation \(\rho_{q}\) is equal to \(K_{pq}\) .
\end{enumerate}

\(f\) ) Prove that the degree of \(\lambda_{\mathbf{p}}\) is \(\frac{n!}{\mathbf{p}^{\prime}!}\Delta (\mathbf{p}^{\prime})\) (expand the polynomial \(\Delta (\mathbf{X})(\sum \mathrm{X}_i)^n\) ).

¶ 25) Consider the representation \(\pi_{\mathbf{n}}\) of \(\mathfrak{S}_n\) associated with the action of \(\mathfrak{S}_n\) on the set \(\mathbf{n} = [1, n]\) (Exercise 6). Let \(k\) be a positive integer and \(\chi_k\) the character of the representation \(\wedge^k \pi_{\mathbf{n}}\) .

\begin{enumerate}
\def\labelenumi{\alph{enumi})}
\item
  Let \(\sigma \in \mathfrak{S}_n\) . Prove the formula \(\chi_k(\sigma) = \sum \varepsilon (\sigma |S)\) , where the sum is taken over the set of subsets \(S\) of \(\mathbf{n}\) with \(k\) elements such that \(\sigma (S) = S\) .
\item
  Prove that \(\wedge^{k}\pi_{n}\) is the irreducible representation of \(\mathfrak{S}_n\) associated with the element \((n-k,1,\ldots,1)\) of \(\mathcal{S}_n\) (calculate the character of this representation using the Frobenius formula, observing that the only terms of the sum
\end{enumerate}

\[
\Delta (\mathbf {X}) f ^ {\boldsymbol {\alpha}} (\mathbf {X}) = \sum_ {\sigma \in \mathfrak {S} _ {n}} \varepsilon (\sigma) X _ {1} ^ {\sigma (n) - 1} \dots X _ {n} ^ {\sigma (1) - 1} f ^ {\boldsymbol {\alpha}} (\mathbf {X})
\]

for which the coefficient of \(X^{p'}\) is nonzero correspond to the permutations \(\sigma\) that satisfy \(\sigma(i) \leqslant i + 1\) for every i and \(\sigma(i) = i\) for \(1 \leqslant i < n - k\) .

¶ 26) Let q be a power of a prime number and F a finite field with q elements; take G to be the group \(\mathbf{GL}(2, \mathbf{F})\) .

\begin{enumerate}
\def\labelenumi{\alph{enumi})}
\item
  Describe the conjugacy classes of G (distinguish four types of classes, containing, respectively, 1, \(q^{2}-1\) , \(q^{2}+q\) , and \(q^{2}-q\) elements).
\item
  The group G acts on the projective line P over F; consider the representation \(\pi_{P}\) of G (Exercise 6). For any group homomorphism \(\lambda : F^{*} \to K^{*}\) , denote by \(\delta_{\lambda}\) the representation of degree 1 associated with the homomorphism \(\lambda \circ det\) from G to \(K^{*}\) , and set \(\pi_{\lambda} = \pi_{P} \otimes \delta_{\lambda}\) . Prove that \(\pi_{\lambda}\) is irreducible, and calculate its character.
\item
  Let B be the subgroup of G consisting of the upper triangular matrices. For \(\lambda,\mu\) in \(\mathrm{Hom}(\mathrm{F}^{*},\mathrm{K}^{*})\) , denote by \(\beta_{\lambda,\mu}\) the representation of B of degree 1 defined by \(\beta_{\lambda,\mu}((a_{ij}))=\lambda(a_{11})\mu(a_{22})\) , and set \(\pi_{\lambda,\mu}=\mathrm{Ind}_{\mathrm{B}}^{\mathrm{G}}(\beta_{\lambda,\mu})\) . Prove that \(\pi_{\lambda,\mu}\) is irreducible for \(\lambda\neq\mu\) , while \(\pi_{\lambda,\lambda}\) is isomorphic to \(\pi_{P}\oplus\delta_{\lambda}\) (apply Exercise 8). The representations \(\pi_{\lambda,\mu}\) and \(\pi_{\lambda',\mu'}\) are isomorphic if and only if the subsets \(\{\lambda,\mu\}\) and \(\{\lambda',\mu'\}\) of \(\mathrm{Hom}(\mathrm{F}^{*},\mathrm{K}^{*})\) are equal.
\end{enumerate}

\(d\) ) Let \(\mathrm{F}^{\prime}\) be an extension of F of degree 2; choosing a basis of \(\mathrm{F}'\) over F allows one to identify \(\mathrm{F}^{\prime *}\) with a subgroup of \(\mathrm{G} = \mathrm{Aut}_{\mathrm{F}}(\mathrm{F}')\) . Let \(\varphi \in \mathrm{Hom}(\mathrm{F}^{\prime *},\mathrm{K}^{*})\) ; calculate \(\mathrm{Ind}_{\mathrm{F}^{\prime *}}^{\mathrm{G}}(\varphi)\) , and prove that it is equal to \(\mathrm{Ind}_{\mathrm{F}^{\prime *}}^{\mathrm{G}}(\varphi^q)\) . Denote by \(\lambda\) the restriction of \(\varphi\) to \(\mathrm{F}^*\) , and set \(\chi_{\varphi} = \chi_{\pi_{\mathrm{P}}\otimes \pi_{\lambda ,1}} - \chi_{\pi_{\lambda ,1}} - \mathrm{Ind}_{\mathrm{F}^{\prime *}}^{\mathrm{G}}(\varphi)\) . Prove that if \(\varphi \neq \varphi^q\) , then we have \(\langle \chi_{\varphi},\chi_{\varphi}\rangle = 1\) and \(\chi_{\varphi}(1) = q - 1\) , so that \(\chi_{\varphi}\) is the character of an irreducible representation \(\pi_{\varphi}\) of dimension \(q - 1\) . Prove that we thus obtain \(\frac{1}{2} q(q - 1)\) nonisomorphic irreducible representations.

\begin{enumerate}
\def\labelenumi{\alph{enumi})}
\setcounter{enumi}{4}
\tightlist
\item
  Prove that the irreducible representations \(\pi_{\lambda}\) for \(\lambda\in\mathrm{Hom}(\mathrm{F}^{*},\mathrm{K}^{*})\) , \(\pi_{\lambda,\mu}\) for a subset \(\{\lambda,\mu\}\) of \(\mathrm{Hom}(\mathrm{F}^{*},\mathrm{K}^{*})\) with two elements, and \(\pi_{\varphi}\) for \(\varphi\in\mathrm{Hom}(\mathrm{F}^{\prime*},\mathrm{K}^{*})\) with \(\varphi\neq\varphi^{q}\) are pairwise nonisomorphic and that every irreducible representation of G is isomorphic to one of them.
\end{enumerate}

\begin{enumerate}
\def\labelenumi{\arabic{enumi})}
\setcounter{enumi}{26}
\tightlist
\item
  Let F be an algebraically closed field of characteristic p \textgreater{} 0 and G be a finite group.
\end{enumerate}

\begin{enumerate}
\def\labelenumi{\alph{enumi})}
\item
  Let \(x \in G\) . Prove that there exists a unique decomposition \(x = x_{u}x_{s}\) , where the order of \(x_{u}\) is a power of p, that of \(x_{s}\) is prime to p, and \(x_{u}\) and \(x_{s}\) commute. The elements \(x_{u}\) and \(x_{s}\) are powers of x.
\item
  Let \(\pi\) be a linear representation of \(G\) over \(F\) of finite degree. Prove that we have \(\chi_{\pi}(x) = \chi_{\pi}(x_s)\) for every \(x \in G\) .
\item
  We say that an element x of G is p-regular if its order is prime to p; we denote by \(G^{reg}\) the set of p-regular elements of G and by \(\mathcal{Z}_{\mathrm{F}}(\mathrm{G}^{\mathrm{reg}})\) the space of central functions from \(G^{reg}\) to F. From the homomorphism \(\mathrm{R}_{\mathrm{F}}(\mathrm{G}) \to \mathcal{Z}_{\mathrm{F}}(\mathrm{G}^{\mathrm{reg}})\) that sends the class of a representation \(\pi\) to the restriction of \(\chi_{\pi}\) to \(G^{reg}\) , we deduce a ring homomorphism \(\psi : F \otimes_{\mathbf{Z}} R_{\mathrm{F}}(G) \to \mathcal{Z}_{\mathrm{F}}(\mathrm{G}^{\mathrm{reg}})\) . Prove that \(\psi\) is injective (use b) and the corollary of VIII, p.~384).
\item
  Let x be a p-regular element of G, \(Z(x)\) its commutant, P a Sylow p-subgroup of \(Z(x)\) , and \(H_{x}\) the subgroup of G generated by P and x. Prove that \(H_{x}\) is the direct product of P and the subgroup of G generated by x and that its conjugacy class is well determined by the conjugacy class of x.
\item
  Construct representations \(\sigma_{1},\ldots,\sigma_{m}\) of \(H_{x}\) of degree 1 and elements \(a_{1},\ldots,a_{m}\) of F such that the function \(\sum a_{i}\chi_{\sigma_{i}}\) takes value 1 at x and 0 at the other powers of x (first construct representations of the subgroup generated by x). Let \(\rho_{i}=\operatorname{Ind}_{H_{x}}^{G}\pi_{i}\) for \(i=1,\ldots,m\) , and let \(f=\sum_{i}a_{i}\chi_{\rho_{i}}\) . Prove that \(f(x)\) is equal to \((\mathrm{Z}(x):\mathrm{H}_{x})1_{\mathrm{F}}\) , which is nonzero in F, and that \(f(y)\) is zero for every p-regular element of G that is not conjugate to x.
\end{enumerate}

\(f\) ) Deduce from \(e\) ) that the homomorphism \(\psi : \mathrm{F} \otimes_{\mathbf{Z}} \mathrm{R}_{\mathrm{F}}(\mathrm{G}) \to \mathcal{Z}_{\mathrm{F}}(\mathrm{G}^{\mathrm{reg}})\) is bijective. In particular, the number of isomorphism classes of simple F{[}G{]}-modules is equal to the number of conjugacy classes of \(p\) -regular elements of G.

\begin{enumerate}
\def\labelenumi{\arabic{enumi})}
\setcounter{enumi}{27}
\tightlist
\item
  Keep the assumptions of Exercise 27. Let \(m\) be the l.c.m. of the orders of the \(p\) -regular elements of \(G\) . Let \(F_0\) be the subfield of \(F\) generated by the \(m\) -th roots of unity; it is a finite field.
\end{enumerate}

\begin{enumerate}
\def\labelenumi{\alph{enumi})}
\item
  Prove that the character of any representation of G over F of finite degree takes all its values in \(F_{0}\) .
\item
  Prove that the quotient ring of \(F_{0}[G]\) by its radical is a product of matrix algebras over \(F_{0}\) (use a) and Wedderburn's theorem). Deduce that the canonical homomorphism \(\mathrm{R}_{\mathrm{F}_{0}}(\mathrm{G}) \to \mathrm{R}_{\mathrm{F}}(\mathrm{G})\) is bijective.
\end{enumerate}

\begin{enumerate}
\def\labelenumi{\arabic{enumi})}
\setcounter{enumi}{28}
\tightlist
\item
  Keep the assumptions and notation of Exercises 27 and 28. Choose a primitive \(m\) -th root of unity \(\zeta\) . Let \(\mathcal{O}_m\) be the ring \(\mathbf{Z}[\mathrm{X}] / (\Phi_m)\) (VIII, p.~415); we denote the class of X in \(\mathcal{O}_m\) by \(x\) . Let \(\varphi : \mathcal{O}_m \to \mathrm{F}\) be the homomorphism that sends \(x\) to \(\zeta\) . It induces an isomorphism from the subgroup of \(\mathcal{O}_m\) generated by \(x\) to the group \(\mu_m(\mathrm{F})\) ; we denote the inverse isomorphism by \(\tau\) .
\end{enumerate}

\begin{enumerate}
\def\labelenumi{\alph{enumi})}
\item
  Let \(\pi\) be a linear representation of G over F, of finite dimension \(d\) . For \(g \in \mathrm{G}\) , let \(\mathrm{P}_g(\mathrm{X}) = \prod_{i=1}^d (\mathrm{X} - \lambda_i)\) be the characteristic polynomial of \(\pi(g)\) ; denote by \(\beta_{\pi}(g)\) the element \(\sum_{i} \tau(\lambda_i)\) of \(\mathcal{O}_m\) . The resulting function \(\beta_{\pi}: \mathrm{G} \to \mathcal{O}_m\) is called the Brauer character of \(\pi\) ; we have \(\varphi \circ \beta_{\pi} = \chi_{\pi}\) and \(\beta_{\pi}(x) = \beta_{\pi}(x_s)\) for every \(x\) , so that \(\beta_{\pi}\) is determined by its restriction to \(\mathrm{G}^{\mathrm{reg}}\) .
\item
  Let \(\widehat{\mathbf{G}}\) be the set of isomorphism classes of simple F{[}G{]}-modules. Prove that the functions \((\beta_{\lambda})_{\lambda \in \widehat{\mathbf{G}}}\) are linearly independent over \(\mathbf{Z}\) . (If there exists a nontrivial relation \(\sum n_{\lambda} \beta_{\lambda} = 0\) , then we may assume that one of the \(n_{\lambda}\) is prime to \(p\) and apply \(\varphi\) to arrive at a contradiction using Exercise 27, f.)
\item
  Let \(\mathcal{B}\) be the set of functions from \(\mathrm{G}^{\mathrm{reg}}\) to \(\mathcal{O}_m\) ; the homomorphism \(\mathrm{R}_{\mathrm{F}}(\mathrm{G}) \to \mathcal{B}\) that sends the class of \(\pi\) to \(\beta_{\pi}\) is injective. Deduce that two semisimple F{[}G{]}-modules are isomorphic if and only if their Brauer characters are equal.
\end{enumerate}

\begin{enumerate}
\def\labelenumi{\arabic{enumi})}
\setcounter{enumi}{29}
\tightlist
\item
  Let \(\mathrm{F}\) be an algebraically closed field of characteristic \(p > 0\) , let \(\mathbf{F}_p\) be the prime subfield of \(\mathrm{F}\) , and let \(\mathrm{G} = \mathbf{SL}_2(\mathbf{F}_p)\) . Denote by \(\mathrm{V}\) the F-vector space \(\mathrm{F} \otimes_{\mathbf{F}_p} \mathbf{F}_p^2\) endowed with the action of \(\mathrm{G}\) deduced from the action on \(\mathbf{F}_p^2\) .
\end{enumerate}

\begin{enumerate}
\def\labelenumi{\alph{enumi})}
\item
  Prove that the F{[}G{]}-module \(\mathbf{S}^{d}(V)\) is simple for \(0 \leqslant d < p\) . (Let W be a nonzero F{[}G{]}-submodule of \(\mathbf{S}^{d}(V)\) , and let \((e, f)\) be a basis of V. Prove that \(e^{d}\) belongs to W using the matrix \(\begin{pmatrix}1 & 1 \\ 0 & 1\end{pmatrix}\) , and then that W is equal to \(\mathbf{S}^{d}(V)\) using the matrix \(\begin{pmatrix}1 & 0 \\ 1 & 1\end{pmatrix}\) .) b) Prove that the F{[}G{]}-module \(\mathbf{S}^{p}(V)\) is not simple (consider the image of V by the mapping \(v \mapsto v^{p}\) ).
\setcounter{enumi}{2}
\item
  For \(p \geqslant 7\) , prove that the dimension of \(\mathbf{S}^{p-3}(\mathrm{V})\) does not divide the order of G. d) Prove that every simple F{[}G{]}-module is isomorphic to one of the modules \(\mathbf{S}^{d}(\mathrm{V})\) for \(0 \leqslant d < p\) (use Exercise 27, f)).
\end{enumerate}

\appendix
\renewcommand{\thechapter}{\arabic{chapter}}
